\chapter{引言：3D表示的演进与动态物理的挑战}

计算机图形学与计算机视觉的交叉领域正在快速发展。长期以来，三维场景的数字化重建与交互式仿真被视为两个相对独立的任务：前者致力于捕捉光影与几何的视觉真实感，后者则专注于力学属性与运动规律的物理真实感。然而，随着元宇宙、数字孪生以及VR/AR应用需求的爆发，学术界与工业界迫切需要一种既能实现照片级真实感渲染，又能遵循物理定律进行交互的统一表达形式。

在这一背景下，3D高斯泼溅（3D Gaussian Splatting, 3DGS）\cite{wu2025survey} 技术的横空出世，为解决这一难题提供了全新的思路。不同于传统的基于网格的显式几何表示，也不同于NeRF的隐式连续表示，3DGS通过离散的、各向异性的高斯椭球体集合来表达场景，兼具了显式表示的高效渲染能力与隐式表示的优化灵活性。

然而，原始的3DGS主要针对静态场景设计。如何将其扩展至动态场景（4D重建），以及如何赋予这些视觉基元物理属性，成为了当前研究的最前沿。本综述报告将讨论该领域的三个里程碑式工作：

\begin{enumerate}
    \item Gaussian-Flow \cite{lin2024gaussianflow_cvpr}：解决了动态场景的高效重建问题，提出了双域变形模型（Dual-Domain Deformation Model）。
    \item PhysGaussian \cite{xie2024physgaussian_cvpr}：实现了物理仿真与视觉渲染的统一，提出了“所见即所仿”的理念。
    \item PhysDreamer \cite{li2024physdreamer_arxiv}：探索了如何利用视频生成模型的先验知识，从静态观测中“蒸馏”出物体的物理材质属性，实现生成式的交互仿真。
\end{enumerate}

\section{从静态到动态：表示方法的困境}

在3DGS出现之前，动态场景的重建主要依赖于动态NeRF（Dynamic NeRF）。虽然动态NeRF能够生成高质量的自由视点视频，但其隐式神经网络的特性导致了极高的计算成本，难以满足实时交互的需求。

随着3DGS的引入，研究人员迅速尝试将其应用于动态场景。早期的尝试主要分为两类：

\begin{enumerate}
    \item 逐帧优化（Per-frame Optimization）：为视频的每一帧训练一个独立的3DGS模型。这种方法虽然能保留极高的细节，但存储空间随视频时长线性增长。对于一段几分钟的视频，其模型大小可能达到数十GB，且训练时间极长，缺乏帧间的时序一致性 \cite{lin2024gaussianflow_cvpr}。
    \item 隐式变形场（Implicit Deformation Fields）：利用一个坐标及时的多层感知机（MLP）来预测每个高斯点在时间t的位移。虽然这种方法节省了存储空间，但引入MLP使得推理过程变慢，因为渲染每一帧都需要查询数百万次神经网络，严重拖慢了渲染FPS \cite{lin2024gaussianflow_cvpr}。
\end{enumerate}

因此，如何在保持3DGS实时渲染优势的同时，高效地通过紧凑的参数来表达复杂的时空动态，成为了Gaussian-Flow所要解决的核心问题。

\section{物理与视觉的鸿沟}

另一个长期存在的挑战是视觉几何与物理几何的分离。在传统的电影特效或游戏开发管线中，视觉艺术家制作高精度的三角网格（Visual Mesh）用于渲染，而物理引擎则需要一套简化的、通常是四面体网格（Simulation Mesh）或体素代理来进行碰撞和形变计算。

这种分离导致了繁琐的“绑定”（Binding）过程，且常常出现视觉模型与物理模型不匹配的伪影（Artifacts）。PhysGaussian 提出了一种观点：如果我们的渲染基元（3D高斯）本身就是物理粒子，那么这种鸿沟将不复存在 \cite{xie2024physgaussian_cvpr}。这种统一不仅简化了管线，还使得物理仿真能够直接驱动高保真的视觉效果。

\section{数据匮乏与生成式先验}

即使拥有了完美的仿真引擎，物理参数（如杨氏模量、泊松比、密度）的获取依然是一个巨大的难题。对于现实世界中捕捉到的物体（例如一段视频中的花朵），我们很难通过非接触手段准确测量其物理属性。

PhysDreamer 和 Generative Image Dynamics \cite{li2024generative_cvpr} 代表了生成式AI与物理仿真结合的新方向。它们假设：在大规模视频数据上训练的视频生成模型（如Sora）已经隐式地学习到了世界的物理规律。通过精心设计的“蒸馏”算法，我们可以将这些隐式先验转化为显式的物理参数，从而实现对未知物体的逼真交互仿真 \cite{li2024physdreamer_arxiv}。
